% !TeX root = ../数模通用模板.tex
\section{问题背景与问题重述}
\subsection{问题背景}
分布式光伏为社区提供了就地获取电能的途径，但其电力供给的时间分布与居民用电需求并不天然一致：发电充裕时可能出现剩余电量，而负载集中或光伏回落时又需要依赖外部电网。配置储能电池可以将不同时段的电能供需连接起来，却也引入了充放电损耗、功率上限和安全储电范围等限制。题目给出的储能设备最大容量为$12000$\,kWh，最大充放电功率为$5000$\,kW，运行过程中储电量须保持在$1200$--$10800$\,kWh之间，充放电效率为$90\%$，并给定了评价期初始储电量。这些参数使得相邻时段的调度决策相互关联，购电、充电、放电和剩余电量必须同时满足能量平衡与设备运行约束。

题目同时提供了不同层次的外部信息：附件1给出典型日的分时电价、社区负载和光伏预测，附件2给出全年按时段记录的负载与光伏实际功率，附件3给出每天多个时刻发布的未来24小时光伏预报，附件4进一步给出随日期和时段波动的外网电价。因而，调度者面对的并不是一次性已知的全天数据，而是随着时间推进逐步获得信息的决策过程。若计划购电量不足以覆盖实际负载，就需要以交易时刻电价的5倍进行紧急购电；若频繁根据新预报调整计划，又会产生计划与调整电量之间的违约或超额购电费用。

因此，微网运营的关键并非单纯提高某一时段的自给率，而是在持续满足负载的前提下，协调购电时机、储能状态与预测风险，控制整个运行周期的购电支出。现实调度还存在“计划先制定、信息后到达”的矛盾。购电计划偏乐观可能引发高价紧急购电，过度保守则可能为未被利用的电量支付费用；中途调整购电安排有助于修正计划，但调整本身也有代价。因此，设计一个“预测—计划—修正—执行—结算”一体化的调度分析模型，对建设兼具经济性与可靠性的微网发电系统具有重要意义。


\subsection{问题重述}
本文以光伏电源、社区负载、储能电池及外网组成的系统为对象，以十分钟时段的购电量和储能动作作为主要决策变量，将原题四问归纳为以下递进任务。四个问题由理想化的典型日逐步推广到全年实测数据、滚动预测和波动电价，分别考察确定性调度、预测误差下的日前计划、信息更新下的动态调整以及价格与供需同时变化时的综合决策。

\textbf{问题一}要求在每天0:00利用典型日的电价、社区负载和给定的光伏功率预测，确定全天各十分钟时段的计划购电量以及储能设备的充放电安排。微网提供的电能不得低于社区负载，并要求储能设备在0:00和24:00的储电量相同。我们建立两层混合整数线性规划模型，在日初日末储电量一致的条件下，先求最低购电费，再在小幅费用让步范围内改善电池充放电的平稳性。购电量作为题目要求的结果进行统计，不再单独最小化，并按题目给定的表1、表2格式汇总指定时段和全天结果。

\textbf{问题二}要求在每天0:00制定全天计划，而当天的社区负载和光伏实际功率将在运行过程中逐步实现。我们先依据历史负载和光伏数据构建预测模型，根据预测得到的每日负载和发电数据制定计划购电量与储能安排；执行时，若计划供电低于实际负载，则需要在相应时段按交易电价5倍进行紧急购电。除紧急购电费用外，其余时段按照0:00锁定的计划购电量结算。该问题需要利用附件1和附件2完成2025年全年滚动评价，并按题目给定格式报告指定日期的计划购电、充放电、全天费用及紧急购电区间。

\textbf{问题三}引入每日0、6、12、18时发布的未来24小时整点光伏预报，研究是否以及如何利用新信息修正尚未执行的购电计划。根据0:00预报制定的计划与后续调整之间存在不对称费用：计划购电量高于调整购电量的部分按交易时刻电价的50\%计收违约费用，调整购电量超过计划购电量的部分按交易时刻电价的1.5倍计收；此外还要计入正常计划购电费用和可能发生的紧急购电费用。该问题重点衡量预报提前量、未来更新机会与调整成本之间的关系，判断是否有必要使用6:00、12:00和18:00的额外预报，并评价这些更新信息的实际决策价值。

\textbf{问题四}将外网的固定日内价格扩展为随日期和时段变化的实时电价，在明确价格可见性的前提下，分别重新研究无额外光伏预报和有日内预报的两类策略。我们利用附件2中的社区负载与光伏实际功率、附件3中的光伏预报以及附件4中的全年电价数据，重新计算问题二和问题三，比较波动价格、供需预测误差及信息更新时间共同作用下的购电安排、储能运行和总费用变化，并将两类策略的完整结果分别保存为题目要求的文件。
